\input{./lib/hw-setup/tex/HWSetup}
\input{./lib/hw-setup/tex/EngBindings}

%
% Homework Details
%   - Title
%   - Subtitle
%   - Due date
%   - Due time
%   - Course
%   - Section/Time
%   - Instructor
%   - Author
%

\hwkTitle{HW02}
\hwkSubTitle{Propulsive Glide Travel from Clavius to Tycho}
\hwkDueDate{2026-02-27}
\hwkDueTime{11:59:00}
\hwkClass{ENAE488X - 0101}
\hwkClassTime{17:00:00}
\hwkInstructor{Dr. Akin}
\hwkAuthor{Vai Srivastava}
\hwkCompletionDate{\today}

\begin{document}

\maketitle

\pagebreak

\section*{Overview} \label{sec:overview}

In the \textit{2001: A Space Odyssey} video clip of propulsive glide shown in \href{https://github.com/vaisriv/uni/blob/main/26s/ENAE488X/docs/lectures/lec02.pdf}{Lecture 2}, the vehicle was traveling from Clavius to Tycho, a distance of \qty{470}{\km}. For the time being, assume both departure and destination are at the same elevation, that the Moon is essentially flat, and that the propulsion systems use \ce{LOX}/\ce{LH2} propellants with an exhaust velocity \( v_{e} = \qty{4500}{\m\per\s} \). Lunar gravity is \qty{1.622}{\m\per\s\squared}.

\begin{hwkProblem*}{1}{} \label{hwk:p01}

	Assuming this mission was an optimum propulsive glide,
	\begin{enumerate}[label=\alph*.]
		\item \label{hwk:p01a} What horizontal velocity would the vehicle have?
		\item \label{hwk:p01b} How long would the flight take (in minutes)?
		\item \label{hwk:p01c} What would the total \( \Delta v \) be for this trip?
		\item \label{hwk:p01d} If the vehicle has an empty mass of \qty{4000}{\kg}, how much propellant would be required for this trip?
	\end{enumerate}

	\hwkSol{} \label{hwk:s01}

	Given the following:
	\begin{align*}
		d &= \qty{470}{\km} = \qty{470e3}{\m} \\
		g &= \qty{1.622}{\m\per\s\squared}
	\end{align*}

	\hwkPart{} \label{hwk:s01a}

	Using:
	\begin{align*}
		V_{\text{opt}} &= \sqrt{\frac{gd}{2}} \\
		\, &= \qty{617.4}{\m\per\s}
	\end{align*}

	We have:
	\[
		\boxed{\text{Horizontal Velocity} = \qty{617.4}{\m\per\s}}
	\]

	\hwkPart{} \label{hwk:s01b}

	Using:
	\begin{align*}
		t_{\text{flight}} &= \frac{d}{V} \\
		\, &= \qty{761.3}{\s}
	\end{align*}

	We have:
	\[
		\boxed{\text{Flight Time} = \qty{12.69}{\min}}
	\]

	\hwkPart{} \label{hwk:s01c}

	Using:
	\begin{align*}
		\Delta V_{\text{total}} &= \Delta V_{h} + \Delta V_{v} \\
		\, &= 2 V + g t_{\text{flight}} \\
		\, &= \qty{2470}{\m\per\s}
	\end{align*}

	We have:
	\[
		\boxed{\text{Total } \Delta V = \qty{2470}{\m\per\s}}
	\]

	\hwkPart{} \label{hwk:s01d}

	Using:
	\begin{align*}
		\Delta V_{\text{total}} &= V_{e} \log{\frac{m_{o}}{m_{f}}} \\
		\implies \frac{m_{f} + m_{pr}}{m_{f}} &= \exp{\frac{\Delta V_{\text{total}}}{V_{e}}} \\
		\implies &= m_{pr} = m_{f} \exp{\frac{\Delta V_{\text{total}}}{V_{e}}} - m_{f} \\
		\, &= \qty{2925}{\kg}
	\end{align*}

	We have:
	\[
		\boxed{\text{Propellant Mass} = \qty{2925}{\kg}}
	\]

\end{hwkProblem*}

\begin{hwkProblem*}{2}{} \label{hwk:p02}

	If the vehicle made the trip using a single ballistic hop,
	\begin{enumerate}[label=\alph*.]
		\item \label{hwk:p02a} What \( \Delta v \) would be required for the initial launch?
		\item \label{hwk:p02b} How long would the flight take (in minutes)?
		\item \label{hwk:p02c} What would be the maximum height the vehicle reaches between departure and destination?
		\item \label{hwk:p02d} What is the total \( \Delta v \) for this trip?
		\item \label{hwk:p02e} If the vehicle has an empty mass of \qty{4000}{\kg}, how much propellant would be required for this trip?
	\end{enumerate}

	\hwkSol{} \label{hwk:s02}

	Given the following:
	\begin{align*}
		\gamma &= \qty{45}{\deg}
	\end{align*}

	\hwkPart{} \label{hwk:s02a}

	Using:
	\begin{align*}
		\Delta V_{\text{launch}} &= \sqrt{g d} \\
		\, &= \qty{873.1}{\m\per\s}
	\end{align*}

	We have:
	\[
		\boxed{\text{Launch Velocity} = \qty{873.1}{\m\per\s}}
	\]

	\hwkPart{} \label{hwk:s02b}

	Using:
	\begin{align*}
		t_{\text{flight}} &= \frac{2 V_{v}}{g} \\
		\, &= \frac{\Delta V_{\text{total}}}{\sqrt{2} g} \\
		\, &= \qty{761.7}{\s}
	\end{align*}

	We have:
	\[
		\boxed{\text{Flight Time} = \qty{12.69}{\min}}
	\]

	\hwkPart{} \label{hwk:s02c}

	Using:
	\begin{align*}
		h_{\text{max}} &= \frac{d}{4} \\
		\, &= \qty{117.5}{\km}
	\end{align*}

	We have:
	\[
		\boxed{\text{Maximum Height} = \qty{117.5}{\km}}
	\]

	\hwkPart{} \label{hwk:s02d}

	Using:
	\begin{align*}
		\Delta V_{\text{total}} &= 2 V \\
		\, &= \qty{1746}{\m\per\s}
	\end{align*}

	We have:
	\[
		\boxed{\text{Total } \Delta V = \qty{1746}{\m\per\s}}
	\]

	\hwkPart{} \label{hwk:s02e}

	Using:
	\begin{align*}
		m_{pr} = m_{f} \left(\exp{\frac{\Delta V_{\text{total}}}{V_{e}}} - 1\right) \\
		\, &= \qty{1896}{\kg}
	\end{align*}

	We have:
	\[
		\boxed{\text{Propellant Mass} = \qty{1896}{\kg}}
	\]

\end{hwkProblem*}

\begin{hwkProblem*}{3}{} \label{hwk:p03}

	The Moon is actually a sphere with a radius of \qty{1738}{\km}. Repeat the analysis of \hyperref[hwk:p01]{Problem 1} assuming a spherical planet instead of a flat one.
	\begin{enumerate}[label=\alph*.]
		\item \label{hwk:p03a} What horizontal velocity would the vehicle have?
		\item \label{hwk:p03b} How long would the flight take (in minutes)?
		\item \label{hwk:p03c} What would the total \( \Delta v \) be for this trip?
		\item \label{hwk:p03d} If the vehicle has an empty mass of \qty{4000}{\kg}, how much propellant would be required for this trip?
	\end{enumerate}

	\hwkSol{} \label{hwk:s03}

	Given the following:
	\begin{align*}
		r &= \qty{1738}{\km}
	\end{align*}

	\hwkPart{} \label{hwk:s03a}

	Using:
	\begin{align*}
		V_{\text{opt}} &= \sqrt{\frac{gd}{2 - \frac{d}{r}}} \\
		\, &= \qty{663.9}{\m\per\s}
	\end{align*}

	We have:
	\[
		\boxed{\text{Horizontal Velocity} = \qty{663.9}{\m\per\s}}
	\]

	\hwkPart{} \label{hwk:s03b}

	Using:
	\begin{align*}
		t_{\text{flight}} &= \frac{d}{V} \\
		\, &= \qty{707.9}{\s}
	\end{align*}

	We have:
	\[
		\boxed{\text{Flight Time} = \qty{11.80}{\min}}
	\]

	\hwkPart{} \label{hwk:s03c}

	Using:
	\begin{align*}
		\Delta V_{\text{total}} &= 2 \sqrt{2 - \frac{d}{r}} \sqrt{gd} \\
		\, &= \qty{2297}{\m\per\s}
	\end{align*}

	We have:
	\[
		\boxed{\text{Total } \Delta V = \qty{2297}{\m\per\s}}
	\]

	\hwkPart{} \label{hwk:s03d}

	Using:
	\begin{align*}
		m_{pr} = m_{f} \left(\exp{\frac{\Delta V_{\text{total}}}{V_{e}}} - 1\right) \\
		\, &= \qty{2663}{\kg}
	\end{align*}

	We have:
	\[
		\boxed{\text{Propellant Mass} = \qty{2663}{\kg}}
	\]

\end{hwkProblem*}

\begin{hwkProblem*}{4}{} \label{hwk:p04}

	A Mars aircraft uses a jet engine which burns \ce{CO2} of the atmosphere with \ce{H2} fuel carried onboard. This combination has a specific fuel consumption of \qty{9.9e-5}{\per\s}. Since hydrogen is very low density, only \qty{10}{\percent} of the aircraft's weight at takeoff is fuel. It has an \( L/D \) of 20 and files at \qty{200}{\m\per\s}. The gravity on Mars is \qty{3.721}{\m\per\s\squared}. Calculate the range of the aircraft.

	\hwkSol{} \label{hwk:s04}

	Given the following:
	\begin{align*}
		\text{SFC} &= \qty{9.9e-5}{\per\s} \\
		L/D &= 20 \\
		V &= \qty{200}{\m\per\s} \\
		g &= \qty{3.721}{\m\per\s\squared} \\
		W_{f} &= 0.9 W
	\end{align*}

	Using:
	\begin{align*}
		R &= \frac{V \frac{L}{D}}{g \, \text{SFC}} \log{\frac{W_{i}}{W_{f}}} \\
		\, &= \qty{1144}{\km}
	\end{align*}

	We have:
	\[
		\boxed{\text{Range} = \qty{1144}{\km}}
	\]

\end{hwkProblem*}

% \begin{hwkProblem*}{5}{\textit{(ENAE788X only)}} \label{hwk:p05}
%
% 	Repeat \hyperref[hwk:p02]{Problem 2} if the landing site at Tycho is \qty{3000}{\m} lower in elevation than the launch point at Clavius. You should still assume a flat planet, but you will have to find a solution numerically.
% 	\begin{enumerate}[label=\alph*.]
% 		\item \label{hwk:p05a} What \( \Delta v \) would be required for the initial launch?
% 		\item \label{hwk:p05b} What \( \Delta v \) would be required for the landing?
% 		\item \label{hwk:p05c} How long would the flight take (in minutes)?
% 		\item \label{hwk:p05d} What would be the maximum height the vehicle reaches between departure and destination?
% 		\item \label{hwk:p05e} What is the total \( \Delta v \) for this trip?
% 		\item \label{hwk:p05f} If the vehicle has an empty mass of \qty{4000}{\kg}, how much propellant would be required for this trip?
% 	\end{enumerate}
%
% 	\hwkSol{} \label{hwk:s05}
%
% 	\hwkPart{} \label{hwk:s05a}
%
% 	Answer
%
% 	\hwkPart{} \label{hwk:s05b}
%
% 	Answer
%
% 	\hwkPart{} \label{hwk:s05c}
%
% 	Answer
%
% 	\hwkPart{} \label{hwk:s05d}
%
% 	Answer
%
% 	\hwkPart{} \label{hwk:s05e}
%
% 	Answer
%
% 	\hwkPart{} \label{hwk:s05f}
%
% 	Answer
%
% \end{hwkProblem*}

\end{document}
